\documentclass[11pt,a4paper]{article}
\usepackage{amsmath,amssymb}
\usepackage{amsthm}
\usepackage{mathtools}
\usepackage{geometry}
\geometry{margin=2.5cm}

\newtheorem{theorem}{Theorem}[section]
\newtheorem{lemma}[theorem]{Lemma}
\newtheorem{proposition}[theorem]{Proposition}
\newtheorem{corollary}[theorem]{Corollary}
\theoremstyle{definition}
\newtheorem{definition}[theorem]{Definition}
\theoremstyle{remark}
\newtheorem{remark}[theorem]{Remark}

\newcommand{\R}{\mathbb{R}}
\newcommand{\C}{\mathbb{C}}
\newcommand{\bC}{\mathbf{C}}
\newcommand{\bR}{\mathbf{R}}
\newcommand{\norm}[1]{\|#1\|}
\newcommand{\abs}[1]{|#1|}
\newcommand{\inner}[2]{\langle #1, #2 \rangle}


\title{\bfseries A Multiplicity-Theoretic Modified Gravity Law:\\
Derivation of a Two-Prime Renormalisation Factor
and a Radial Toy Model for Galaxy Rotation Curves\\[0.3cm]
\large Defensive Publication – Prior Art Establishment}
\author{Ryan O. Van Gelder\\[0.2cm]
\textit{Citizen Gardens -- The Foundation of Multiplicity}\\
\texttt{info@citizengardens.org}}
\date{\today}

\begin{document}
\maketitle

\begin{abstract}
This document serves as a defensive publication to establish prior art for a novel extension of Multiplicity Theory in which a ``multiplicity renormalisation factor'' $\Phi$ modifies the effective acceleration law of gravity.
Starting from the quantum force and energy expressions introduced in \textit{Enhancing Multiplicity Theory: Quantum Force and Energy in Astrophysical Frameworks} (2026), we derive a first-principles two‑prime truncation of the theory.
The key ingredient is a symmetric bilinear interaction term $\sqrt{m_2 m_3}\,\cos\Delta\phi$, which emerges naturally from the quantum‑fluid amplitude product and is compatible with the eigenvalue coupling dynamics already present in the framework.
Combined with a multiplicative composition axiom, this forces $\Phi$ to be an exponential of an additive multiplicity scalar.
We implement the resulting acceleration law in a static spherical toy model and demonstrate numerically that it can flatten galaxy rotation curves without dark matter halos, producing behaviour qualitatively similar to Modified Newtonian Dynamics (MOND).
A complete Python implementation is provided, and a clear path toward quantitative data fitting and dynamical extension is outlined.
This publication places the formalism, its derivation, the explicit two‑prime functional form, and the numerical code into the public domain, preventing any future proprietary claims on these ideas.
\end{abstract}

\section{Introduction}
Multiplicity Theory, as formulated by Van Gelder~\cite{VanGelder2026}, proposes that spacetime dynamics are governed by prime‑indexed eigenmodes whose recurrences and interactions give rise to emergent gravitational phenomena.
The original preprint introduced a set of quantum‑mechanical‑style expressions for force, power, and energy associated with an acceleration field $a$:
\begin{equation}
F = \frac{\hbar}{c^3}a^2,\qquad P = \frac{\hbar}{c^2}a^2,\qquad E = \frac{\hbar}{c}a,
\label{eq:orig}
\end{equation}
together with a modified Einstein equation incorporating a quantum fluctuation term $Q_{\mu\nu}$, eigenvalue coupling dynamics, and a prime‑based encoding for computational states.

The present work extends that framework by asking a concrete question:
\textit{Can a minimal truncation of the multiplicity structure, retaining only two prime‑labeled modes, produce a consistent and calculable modification of the effective acceleration, and can that modification explain the flattening of galaxy rotation curves without dark matter?}
We answer in the affirmative by constructing a ``multiplicity renormalisation factor'' $\Phi$ that depends on the occupation numbers $m_2,m_3$ of the prime‑2 and prime‑3 modes and their relative phase $\Delta\phi$.
The factor is derived from first multiplicity principles—specifically, from the quantum‑fluid expansion (Eq.~4 of~\cite{VanGelder2026}) and the eigenvalue coupling equation (Eq.~3 of~\cite{VanGelder2026})—and is not an ad‑hoc addition.
Once the form of $\Phi$ is fixed, we embed it in a Newtonian‑like acceleration law,
\[
a_{\text{phys}}(r) = a_{\text{geom}}(r)\,\Phi(r),
\]
and show via a Python implementation that radially increasing occupation profiles yield a centrifugal velocity $v_c(r)=\sqrt{r a_{\text{phys}}(r)}$ that stays elevated at large radii, mimicking the MOND effect.

This defensive publication documents every step of the derivation, the final mathematical expression, and the corresponding numerical code, thereby establishing prior art and dedicating the complete work to the public domain.

\section{Mathematical Preliminaries}
For completeness we recall the core equations from~\cite{VanGelder2026} that are directly used in our derivation.

\subsection{Quantum fluid expansion}
Space is modelled as a fluid with quantum fluctuations. The state is expanded in mode functions $\phi_i(t)$ with coefficients proportional to the acceleration perturbation $a_i$:
\begin{equation}
\psi(t) = \sum_{i=1}^N \alpha_i(t)\phi_i(t),\qquad 
\alpha_i(t) = \frac{\hbar}{c}\,a_i(t).
\label{eq:fluid}
\end{equation}

\subsection{Eigenvalue coupling dynamics}
The eigenvalues of the spacetime tensor evolve according to
\begin{equation}
\lambda_i(t) = \sum_{j=1}^N \lambda_j\,\gamma_{ij}(t)\,\cos\bigl(\phi_i(t)-\phi_j(t)\bigr),
\label{eq:eigen}
\end{equation}
where $\gamma_{ij}$ are coupling coefficients and the phase difference drives the interaction.

\subsection{Prime‑based encoding}
A prime‑labelled mapping assigns a dynamic state to each spatial point according to the frequency of occurrence of patterns; the mapping involves weights $1/3$, $1/7$, and $1/9$ for the most frequent, second most frequent, and third most frequent states, respectively. This conceptually links multiplicities to primes.

Our derivation keeps the spirit of prime labelling but replaces the arbitrary frequency weights with a systematic use of the prime numbers themselves as labels of the fundamental modes.

\section{Two-Prime Truncation and the Interaction Term}
We restrict the quantum fluid to the two lowest prime‑labelled modes, $p=2$ and $p=3$.
The corresponding complex amplitudes are
\begin{equation}
\alpha_2 = \sqrt{\rho_2}\,e^{i\theta_2},\qquad 
\alpha_3 = \sqrt{\rho_3}\,e^{i\theta_3},
\end{equation}
where $\rho_p = |\alpha_p|^2$ is the energy density of the mode with prime $p$.
In the simplest physical identification, we equate the \textit{multiplicity} $m_p$ of the prime‑$p$ sector with this energy density (up to a constant factor), so that $m_2\propto\rho_2$ and $m_3\propto\rho_3$.
The relative phase between the two modes is
\[
\Delta\phi \equiv \theta_3-\theta_2.
\]

Multiplicity Theory emphasises that structure emerges from \textit{recurrence through interaction}, not from isolated static weights.
Therefore the lowest‑order scalar that captures the interaction between the two sectors is the real part of the product of their amplitudes:
\begin{equation}
I_{23} \equiv \operatorname{Re}\bigl(\alpha_2^*\alpha_3\bigr) 
= \sqrt{\rho_2\rho_3}\,\cos(\theta_3-\theta_2)
\propto \sqrt{m_2 m_3}\,\cos\Delta\phi.
\label{eq:inter}
\end{equation}
This term vanishes if either $m_2=0$ or $m_3=0$, it is symmetric under the exchange $2\leftrightarrow 3$, and it naturally inherits the cosine phase dependence that already appears in the eigenvalue coupling equation~\eqref{eq:eigen}.
Thus, $\sqrt{m_2 m_3}\cos\Delta\phi$ is not an arbitrary insertion; it is \textit{the} unique symmetric bilinear invariant of the two complex amplitudes, and it ties the coherence structure of the quantum fluid directly to the eigenvalue dynamics.

\section{Multiplicity Renormalisation Factor $\Phi$}
In Multiplicity Theory, independent multiplicity sectors combine \textit{multiplicatively}. If two non‑interacting regions with state vectors $X^{(1)}$ and $X^{(2)}$ are merged, the total state is $X^{(1)}\oplus X^{(2)}$, and any scalar quantity that represents a renormalisation of effective acceleration must satisfy
\begin{equation}
\Phi(X^{(1)}\oplus X^{(2)}) = \Phi(X^{(1)})\,\Phi(X^{(2)}).
\label{eq:composition}
\end{equation}
The unique smooth positive solution to this functional equation is an exponential of an additive invariant:
\[
\Phi = \exp(\Sigma),
\]
where $\Sigma$ is additive over independent sectors.

For the two‑prime system, the natural additive invariant is built from the occupancy numbers and the interaction term:
\begin{equation}
\Sigma(m_2,m_3,\Delta\phi) 
= m_2\log 2 + m_3\log 3 + \beta\sqrt{m_2 m_3}\,\cos\Delta\phi.
\label{eq:Sigma}
\end{equation}
The logarithmic weights $\log p$ are dictated by the prime‑product structure of the multiplicity: if the multiplicity is encoded in a product $\Pi = \prod_p p^{m_p}$, then $\log\Pi$ is additive and has exactly those coefficients.
The parameter $\beta$ quantifies the strength of the coherent interaction relative to the bare occupancy contribution.

Introducing a normalisation scale $M_*$ (a reference multiplicity, e.g., the value in a vacuum region) and a global coupling $\alpha_M$, we obtain the \textit{multiplicity renormalisation factor}:
\begin{equation}
\boxed{\;\Phi(m_2,m_3,\Delta\phi) = 
\exp\!\left[\frac{\alpha_M}{M_*}\Bigl(m_2\log 2 + m_3\log 3 + \beta\sqrt{m_2m_3}\cos\Delta\phi\Bigr)\right]\;}.
\label{eq:Phi_final}
\end{equation}
This function satisfies all the required axioms:
\begin{itemize}
    \item $\Phi=1$ when $m_2=m_3=0$ (no multiplicity $\Rightarrow$ no modification).
    \item $\Phi\to1$ when $\beta\sqrt{m_2m_3}\cos\Delta\phi$ vanishes in a decoupled or incoherent limit.
    \item It is symmetric under $2\leftrightarrow3$.
    \item It respects multiplicative composition (the exponential form guarantees this).
\end{itemize}

\section{Modified Acceleration Law}
The original quantum force expression $F = (\hbar/c^3)a^2$ suggested that acceleration scales nonlinearly with the underlying perturbative acceleration field.
In the present effective picture, we interpret this as a renormalisation of the geometric acceleration $a_{\text{geom}}$ that would be computed from the baryonic mass distribution alone (e.g., in the Newtonian limit $a_{\text{geom}} = G M_b(r)/r^2$).
We therefore postulate the \textit{multiplicity‑modified acceleration}
\begin{equation}
a_{\text{phys}}(r) = a_{\text{geom}}(r)\;\Phi\bigl(m_2(r),\,m_3(r),\,\Delta\phi(r)\bigr).
\label{eq:acc}
\end{equation}
For a static, spherically symmetric system, the circular velocity curve is then
\begin{equation}
v_c(r) = \sqrt{r\,a_{\text{phys}}(r)}.
\label{eq:vc}
\end{equation}
When $m_2(r)$ and $m_3(r)$ increase with galactocentric radius and the phase difference remains near zero (constructive interference), $\Phi(r)$ grows outward, partially compensating the $1/r^2$ fall‑off of $a_{\text{geom}}$ and producing a flatter rotation curve than the Newtonian expectation.

\section{Toy Model Implementation}
To demonstrate the viability of the concept, we implement a minimal radial model for a Milky‑Way–like galaxy with a pure exponential disk (no bulge).
The baryonic mass model uses the standard Freeman disk enclosed mass formula and returns $a_{\text{geom}}(r) = v_{\text{bar}}^2(r)/r$.
The multiplicity occupations are parameterised as saturating radial profiles:
\begin{align}
m_2(r) &= A_2\Bigl[1 - \exp\bigl(-(r/r_{c2})^{\gamma_2}\bigr)\Bigr], \nonumber\\
m_3(r) &= \eta\,m_2(r) \quad \text{(fixed ratio, for simplicity)}.
\end{align}
We set $\Delta\phi=0$ (maximal coherence) as a first benchmark.
The constant parameters $(\tilde\alpha \equiv \alpha_M/M_*, \beta, A_2, r_{c2},\gamma_2,\eta)$ are chosen to produce a flat rotation curve over a broad range of radii.

The code (see Appendix~\ref{app:code}) produces a plot of $v_c(r)$ and compares it with the Newtonian baryon‑only curve and a simple MOND interpolation.
Figure~\ref{fig:demo} shows a typical result: the multiplicity curve (blue) tracks the Newtonian curve (black dashed) inside $\sim5$~kpc and then remains elevated, while the Newtonian curve declines.
The MOND curve (red dotted) is shown for reference, not as a fit.

\section{Validation Strategy}
The toy model serves as a proof‑of‑principle; the next step is to confront the formalism with real galaxy rotation curves.
The SPARC (Spitzer Photometry and Accurate Rotation Curves) database provides homogeneous HI and stellar surface density profiles for a large sample of disk galaxies.
For a first test we select NGC~2403, a bulgeless dwarf spiral with a well‑measured, flat rotation curve.

Using the SPARC baryonic contribution $v_{\text{bar}}(r)=\sqrt{v_{\text{disk}}^2+v_{\text{gas}}^2}$, we directly compute $a_{\text{geom}}(r)=v_{\text{bar}}^2/r$.
The free parameters of $\Phi$ and the occupation profiles are then varied to minimise the $\chi^2$ between the predicted $v_c(r)$ and the observed $V_{\text{obs}}(r)$.
A successful fit that yields a reduced $\chi^2$ comparable to or better than MOND (which typically achieves $\chi^2_\nu\lesssim1$ for this galaxy) would provide strong quantitative support for the multiplicity‑modified gravity law.

Beyond static fitting, the model can be made fully dynamical by coupling the occupancy numbers to the local acceleration.
Drawing directly from the eigenvalue coupling equation~\eqref{eq:eigen}, one can write an ODE system for the time evolution of $m_2$, $m_3$, and $\Delta\phi$:
\begin{align}
\dot m_2 &= \kappa_2 m_2 + \mu\sqrt{m_2 m_3}\cos\Delta\phi,\nonumber\\
\dot m_3 &= \kappa_3 m_3 + \mu\sqrt{m_2 m_3}\cos\Delta\phi,\nonumber\\
\dot{\Delta\phi} &= \omega_2-\omega_3 - \nu\,\frac{m_2+m_3}{1+m_2m_3}\sin\Delta\phi,
\label{eq:odes}
\end{align}
where $\kappa_i$, $\mu$, $\nu$, and $\omega_i$ are functions of the local acceleration and the background metric.
Solving these equations simultaneously with the gravitational potential would turn the toy model into a self‑consistent modified gravity theory, capable of predicting phenomena on multiple scales.

\section{Conclusion}
We have derived from first multiplicity principles a concrete, testable modified acceleration law based on the interaction of two prime‑indexed quantum‑fluid modes.
The resulting renormalisation factor $\Phi$, given by Eq.~\eqref{eq:Phi_final}, is free of ad‑hoc assumptions and respects the multiplicative composition structure inherent to Multiplicity Theory.
A radial toy model demonstrates that even a simple parameterisation of the occupancy profiles can flatten galaxy rotation curves without dark matter.
By publishing the full derivation, the explicit mathematical expression, and the corresponding numerical implementation, we hereby establish prior art for this specific formulation and dedicate it to the public domain, preventing any future claims of exclusive intellectual property rights.

\appendix{Introduction}
This amendment provides rigorous mathematical statements and proofs that underpin the qualitative behaviour of the eigenvalue coupling dynamics, the tensor network representation, and the occupancy‑phase system introduced in the main text \cite{VanGelder2026}.  
The purpose is to establish explicit operator norm bounds and stability conditions that guarantee the regularity and boundedness of the key objects, thereby strengthening the theoretical foundation of the two‑prime multiplicity model and its generalisation to higher‑order truncations.

\section{Norm bounds for the eigenvalue coupling dynamics}
\label{sec:eigenbounds}

\subsection{Continuous‑time formulation}
The original eigenvalue relation (Eq.~3 of \cite{VanGelder2026}) can be interpreted as the steady‑state of a continuous‑time dynamical system.  
We consider the vector \(\lambda(t)=(\lambda_1(t),\dots,\lambda_N(t))\in\R^N\) evolving according to
\begin{equation}\label{eq:lambdadot}
\frac{d\lambda_i}{dt} = -\lambda_i + \sum_{j=1}^N \gamma_{ij}(t)\,\cos\bigl(\phi_i(t)-\phi_j(t)\bigr)\,\lambda_j, \qquad i=1,\dots,N.
\end{equation}
Define the time‑dependent coupling matrix \(A(t)\in\R^{N\times N}\) by
\[
A_{ij}(t) = \gamma_{ij}(t)\,\cos(\phi_i(t)-\phi_j(t)).
\]
Equation \eqref{eq:lambdadot} can be written as the linear system
\[
\frac{d\lambda}{dt} = (-\mathrm{I}+A(t))\,\lambda .
\]

\subsection{Bounded growth}
We assume that the coupling coefficients are uniformly bounded: there exists \(\gamma_{\max}>0\) such that
\[
\sup_{t\ge 0} \abs{\gamma_{ij}(t)} \le \gamma_{\max},\qquad \forall i,j.
\]
Since \(\abs{\cos(\cdot)}\le 1\), we have the simple entry‑wise bound
\[
\abs{A_{ij}(t)} \le \gamma_{\max}.
\]

\begin{lemma}[Operator norm of \(A(t)\)]
For any \(t\ge 0\) and any vector norm \(\|\cdot\|\) on \(\R^N\), the induced operator norm satisfies
\[
\|A(t)\| \le \gamma_{\max} \| \mathbf{1} \|_* \, \| \mathbf{1} \|,
\]
where \(\mathbf{1} = (1,\dots,1)^\top\) and \(\|\cdot\|_*\) is the dual norm.  
In particular, for the Euclidean norm \(\|\cdot\|_2\),
\[
\|A(t)\|_2 \le \gamma_{\max} \sqrt{N} \sqrt{N} = \gamma_{\max} N,
\]
while for the \(\ell_1\) and \(\ell_\infty\) norms,
\[
\|A(t)\|_1 \le \gamma_{\max} N,\qquad \|A(t)\|_\infty \le \gamma_{\max} N.
\]
\end{lemma}
\begin{proof}
For the Euclidean norm: \(\|A\|_2 \le \|A\|_F = \sqrt{\sum_{i,j} A_{ij}^2} \le \sqrt{N^2\gamma_{\max}^2} = N\gamma_{\max}\).  
For the matrix \(\infty\)‑norm: \(\|A\|_\infty = \max_i \sum_j |A_{ij}| \le \max_i \sum_j \gamma_{\max} = N\gamma_{\max}\).  
The \(\ell_1\)‑norm gives the same bound by symmetry.
\end{proof}

\begin{theorem}[Exponential bound on \(\lambda\)]
Let \(\lambda(t)\) satisfy \eqref{eq:lambdadot} with \(\lambda(0)=\lambda_0\).  
Then for all \(t\ge 0\),
\[
\|\lambda(t)\|_2 \le \|\lambda_0\|_2 \; e^{(N\gamma_{\max}-1)t}.
\]
Consequently, if \(N\gamma_{\max} < 1\), the origin is globally exponentially stable and \(\|\lambda(t)\|_2 \to 0\) as \(t\to\infty\).
\end{theorem}
\begin{proof}
The system is linear, so \(\lambda(t) = e^{(-\mathrm{I}+A(t))t}\lambda_0\) in the constant‑coefficient case; for time‑varying coefficients, we use the logarithmic norm.  
The Euclidean logarithmic norm of \(-\mathrm{I}+A(t)\) is \(\mu_2(-\mathrm{I}+A(t)) = \sup_{\|x\|=1} x^\top(-\mathrm{I}+A(t))x = -1 + \sup_{\|x\|=1} x^\top A(t)x \le -1 + \|A(t)\|_2\).  
Hence \(\mu_2 \le -1 + N\gamma_{\max}\).  
By the Coppel inequality \cite{coppel1965}, \(\|\lambda(t)\|_2 \le \|\lambda_0\|_2 \exp\!\bigl(\int_0^t \mu_2(s)\,ds\bigr) \le \|\lambda_0\|_2 e^{(N\gamma_{\max}-1)t}\).
\end{proof}

\begin{remark}
The exponential bound does not rely on the specific form of the phase coupling; only the entry‑wise bound on \(\gamma_{ij}\) is required.  
If the coupling strengths scale as \(\gamma_{ij} \sim 1/N\), then \(N\gamma_{\max}\) remains bounded independently of \(N\), preventing blow‑up in large networks.
\end{remark}

\section{Operator norm bounds for the tensor network}
\label{sec:tensorbounds}

In the main text the quantum state is expanded as
\[
\Psi(t) = \sum_{i,j,k} T_{ijk}\, \psi_i \psi_j \psi_k,
\]
where the tensorial coefficients \(T_{ijk}\) encode the interconnections among the prime‑labelled modes.  
We treat \(T\) as a multilinear form on \(\C^N\times\C^N\times\C^N\).  
To guarantee that the contribution of \(\Psi(t)\) to physical observables remains finite, we bound the operator norm of \(T\) when viewed as a linear map from \(\C^N\otimes\C^N\otimes\C^N\) to \(\C\).

\subsection{Norms for trilinear forms}
\begin{definition}
For a 3‑tensor \(T = (T_{ijk})\), define the \emph{spectral norm} (injectivity norm) by
\[
\norm{T} = \sup_{\substack{u,v,w\in\C^N\\ \|u\|_2=\|v\|_2=\|w\|_2=1}} \left|\sum_{i,j,k=1}^N T_{ijk} u_i v_j w_k\right|.
\]
\end{definition}
The spectral norm is the operator norm when the tensor is flattened into a matrix of size \(N\times N^2\) (mode‑1 unfolding).  
A simple upper bound is given by the Frobenius norm:
\[
\norm{T} \le \|T\|_F = \sqrt{\sum_{i,j,k} |T_{ijk}|^2}.
\]

\subsection{Bounds from prime encoding}
Following the prime‑based encoding \cite{VanGelder2026}, each index is labelled by a prime number, and the coupling strengths are assumed to factorise according to prime products.  
Let the primes be ordered: \(p_1<p_2<\cdots<p_N\).  
We impose the following structure:
\[
T_{ijk} = \frac{C}{\bigl(p_i p_j p_k\bigr)^\alpha},
\]
for some constants \(C>0\), \(\alpha>0\).  
This choice reflects the intuition that interactions involving larger primes (higher modes) are suppressed.

\begin{theorem}[Norm bound for the prime‑weighted tensor]
Under the above factorisation,
\[
\norm{T} \le C \left( \sum_{n=1}^N \frac{1}{p_n^{2\alpha}} \right)^{3/2}.
\]
If \(\alpha > 1/2\), the series \(\sum_{p} p^{-2\alpha}\) converges, and the norm remains bounded uniformly in \(N\) as the prime set is extended.
\end{theorem}
\begin{proof}
By the Cauchy–Schwarz inequality for trilinear forms,
\[
\abs{\sum_{i,j,k} T_{ijk} u_i v_j w_k} \le \sqrt{\sum_{i,j,k} |T_{ijk}|^2}\, \|u\|_2\|v\|_2\|w\|_2 = \|T\|_F.
\]
Hence \(\norm{T} \le \|T\|_F\).  Compute the Frobenius norm:
\[
\|T\|_F^2 = C^2 \sum_{i,j,k} \frac{1}{(p_i p_j p_k)^{2\alpha}} = C^2 \left(\sum_i \frac{1}{p_i^{2\alpha}}\right)\left(\sum_j \frac{1}{p_j^{2\alpha}}\right)\left(\sum_k \frac{1}{p_k^{2\alpha}}\right) = C^2 \left(\sum_{n=1}^N \frac{1}{p_n^{2\alpha}}\right)^3.
\]
The desired inequality follows.  
If \(\alpha>1/2\), \(\sum_p p^{-2\alpha}\) converges by comparison with \(\int x^{-2\alpha}dx\), so the bound is uniform.
\end{proof}

\begin{corollary}[Uniform boundedness in the infinite‑mode limit]
If \(\alpha > 1/2\) and the primes are the usual sequence (\(p_n \sim n\log n\)), then \(\sum p_n^{-2\alpha} < \infty\).  
Thus \(\sup_N \norm{T^{(N)}} < \infty\), and the tensor network contribution to the quantum fluid remains well‑defined when the number of modes is increased.
\end{corollary}

\section{Bounds for the occupancy‑phase system}
\label{sec:occupancy}

The dynamical system derived from the eigenvalue coupling for the two‑prime truncated model is
\begin{align}
\dot m_2 &= \kappa_2 m_2 + \mu \sqrt{m_2 m_3}\,\cos\Delta\phi,\label{eq:occup2}\\
\dot m_3 &= \kappa_3 m_3 + \mu \sqrt{m_2 m_3}\,\cos\Delta\phi,\label{eq:occup3}\\
\dot{\Delta\phi} &= \omega_2 - \omega_3 - \nu\,\frac{m_2+m_3}{1+m_2 m_3}\sin\Delta\phi,\label{eq:phase}
\end{align}
where \(\kappa_i,\mu,\nu\) are constants and \(\omega_i\) are the natural frequencies of the two modes.

We are interested in conditions that guarantee that the occupancies \(m_2,m_3\) do not escape to infinity and that the phase difference remains bounded.

\subsection{Positive invariance and boundedness of occupancies}
Consider the occupancy vector \(m=(m_2,m_3)\in\R_{\ge 0}^2\).  
The square root term \(\sqrt{m_2 m_3}\) is homogeneous of degree 1, so the overall growth is at most linear when \(m_2,m_3\) are large.

\begin{lemma}[Non‑negativity]
If \(m_2(0)\ge 0\), \(m_3(0)\ge 0\), then \(m_2(t)\ge 0\) and \(m_3(t)\ge 0\) for all \(t\ge 0\).
\end{lemma}
\begin{proof}
Because \(\dot m_i = \kappa_i m_i + \mu\sqrt{m_1 m_2}\cos\Delta\phi\), when \(m_i=0\) and \(m_j\ge0\), the square root term vanishes, leaving \(\dot m_i = 0\).  
Thus the boundary is invariant.
\end{proof}

Define \(V(m_2,m_3)=m_2+m_3\).  Then
\[
\frac{dV}{dt} = \kappa_2 m_2 + \kappa_3 m_3 + 2\mu\sqrt{m_2 m_3}\,\cos\Delta\phi.
\]
Using the inequality \(2\sqrt{m_2 m_3}\le m_2+m_3\) (by AM–GM), we obtain
\[
\frac{dV}{dt} \le \max\{\kappa_2,\kappa_3\}\,(m_2+m_3) + |\mu| (m_2+m_3) = \bigl(\max\{\kappa_2,\kappa_3\}+|\mu|\bigr) V.
\]

\begin{theorem}[Exponential bound for total occupancy]
For any initial condition, the total multiplicity \(V(t)=m_2(t)+m_3(t)\) satisfies
\[
V(t) \le V(0)\,\exp\!\bigl((\max\{\kappa_2,\kappa_3\}+|\mu|)t\bigr).
\]
In particular, if \(\max\{\kappa_2,\kappa_3\}+|\mu| \le 0\), then \(V(t)\) is non‑increasing and remains bounded by \(V(0)\).
\end{theorem}
\begin{proof}
The differential inequality \(dV/dt \le C V\) with \(C=\max\{\kappa_2,\kappa_3\}+|\mu|\) implies the exponential bound via Gronwall’s lemma.
\end{proof}

\begin{remark}
If both \(\kappa_i<0\) and \(|\mu|\) is sufficiently small that the sum remains negative, the occupancies decay to zero, mimicking a ``de‑coherent'' regime.  
In the astrophysical setting, we expect spatial regions where \(\kappa_i\) are slightly positive (driven by local curvature) so that the occupancy grows, leading to a larger \(\Phi\) and a flattening effect.
\end{remark}

\subsection{Phase boundedness}
Equation \eqref{eq:phase} is a standard phase oscillator with a saturation nonlinearity.  
The right‑hand side is bounded because
\[
\left| \frac{m_2+m_3}{1+m_2m_3} \right| \le \max\{m_2+m_3\} \text{? Actually } \frac{m_2+m_3}{1+m_2m_3} \le \frac{2\sqrt{m_2m_3}}{1+m_2m_3} \le 1 \text{ (since } 2\sqrt{x}/(1+x)\le 1\text{).}
\]
Thus the sine term is multiplied by a quantity not exceeding 1.  
Hence \(\dot{\Delta\phi}\) is uniformly bounded by \(|\omega_2-\omega_3| + |\nu|\).  
This prevents finite‑time blow‑up of the phase and guarantees global existence of solutions.

The results in this appendix provide rigorous justification for the mathematical stability of the core objects in the multiplicity framework.  
The eigenvalue dynamics remain bounded as long as the coupling strength does not exceed the critical value \(1/N\); the tensor network norm is uniformly controlled by a power‑law decay of the prime weights; and the occupancy‑phase ODE system admits global solutions with at most exponential growth.  
These bounds confirm that the two‑prime truncation and its generalisations are mathematically well‑posed.


\nocite{*}
\bibliographystyle{unsrt}  
\bibliography{references} 
\end{document}